\documentclass{article}
\usepackage[utf8]{inputenc}
\usepackage[T1]{fontenc}
\usepackage{geometry}
\usepackage{graphicx}
\usepackage{float}
\usepackage{hyperref}
\usepackage[english,indonesian]{babel}

\title{Optimasi Pemrosesan Template LaTeX Menggunakan Pendekatan Berbasis AI}
\author{Academic Support Agent}
\date{\today}

\begin{document}
\maketitle

\begin{abstract}

\end{abstract}

\chapter{PENDAHULUAN}


\chapter{TINJAUAN PUSTAKA}


% Original content (commented out):
% \documentclass[a4paper]{article}

\usepackage[swedish]{babel}
\usepackage[utf8]{inputenc}
\usepackage{amsmath}
\usepackage{graphicx}
\usepackage[colorinlistoftodos]{todonotes}
\usepackage[]{mcode}
\usepackage[margin=0.8in]{geometry}
\usepackage{mathtools}
\usepackage{float}
\usepackage{tikz}


\title{Title}

\author{Hugo Harlin, Rikard Pihl-Roos}

\date{\today}

\begin{document}

\maketitle

\begin{tikzpicture}[remember picture, overlay]

  \node [anchor=north west, inner sep=30pt]  at (current page.north west)
     {\includegraphics[height=4cm]{uu_logo.png}};
\end{tikzpicture}



\begin{abstract}

I
\end{abstract}

\section{Inledning}



\section{Genomförandet}




\section{Resultat}






\section{Slutsatser}


\section{Appendix}

\begin{lstlisting}



\end{lstlisting}

\end{document}

\end{document}